% Generated draft: requires mathematical and source review.
\documentclass[11pt]{article}
\usepackage[margin=25mm]{geometry}
\usepackage{fontspec}
\setmainfont{Latin Modern Roman}
\usepackage{amsmath,amssymb,mathtools,unicode-math}
\setmathfont{Latin Modern Math}
\usepackage{enumitem,needspace,xurl,hyperref}
\setlength{\parindent}{0pt}
\setlength{\parskip}{5pt}
\setlength{\emergencystretch}{3em}
\newcommand{\N}{\mathbb{N}}
\newcommand{\Z}{\mathbb{Z}}
\newcommand{\Q}{\mathbb{Q}}
\newcommand{\R}{\mathbb{R}}
\newcommand{\C}{\mathbb{C}}
\newcommand{\F}{\mathbb{F}}
\newcommand{\A}{\mathbb{A}}
\newcommand{\PP}{\mathbb{P}}
\newcommand{\E}{\mathbb{E}}
\newcommand{\eps}{\varepsilon}
\newcommand{\dd}{\,\mathrm{d}}
\begin{document}
\section*{UnsolvedMath ID: 701494}
\textbf{Source problem code:} catalogue entry
\subsection{Definitions and mathematical statement}
Let \(p\ne 5\) be a prime and let \(G\) be a finite group. A \(p\)-element of \(G\) is an element whose order is \(p^k\) for some integer \(k\ge 0\), including the identity. Write \(O_p(G)\) for the largest normal subgroup of \(G\) whose order is a power of \(p\). For \(c,d\in G\), use the commutator convention \([c,d]=c^{-1}d^{-1}cd\).

\textbf{Conjecture (attributed in the supplied statement to R. Guralnick and G. Malle).} For every prime \(p\ne 5\), every finite group \(G\), and every conjugacy class
\[
C=\{g^{-1}xg:g\in G\}
\]
of a \(p\)-element \(x\in G\), the implication
\[
\bigl(\text{for all }c,d\in C,\ [c,d]\text{ is a }p\text{-element}\bigr)
\quad\Longrightarrow\quad C\subseteq O_p(G)
\]
holds.

The registry lists V. D. Mazurov as proposer, whereas the statement names R. Guralnick and G. Malle; the supplied record does not explain their respective attribution roles. The registry labels this conjecture open; that label is supplied source metadata, not an independently verified assessment.
\subsection{Short English statement}
For a prime \(p\ne 5\), suppose every commutator of two elements in a conjugacy class of \(p\)-elements of a finite group is again a \(p\)-element. The conjecture asserts that the entire class lies in the largest normal \(p\)-subgroup.
\subsection{Sources}
\begin{itemize}
\item Supplied local UnsolvedMath registry record: permanent ID \(701494\), exact source problem code \texttt{catalogue entry}, title \emph{Kourovka Notebook Problem 18.65}; upstream revision \texttt{d371ee917ac9647fbb06cd6fc529deec6f3dda82}.
\item Supplied source citation: \emph{The Kourovka Notebook (Unsolved Problems in Group Theory)}, 18th issue (2014), Problem 18.65; also described in the record as the 21st issue edition (2026), \texttt{arXiv:1401.0300v48}. The collection description names E. I. Khukhro and V. D. Mazurov as editors.
\item Attribution as supplied: V. D. Mazurov is listed as proposer, with proposed year 2014; the statement attributes the conjecture to R. Guralnick and G. Malle.
\item The supplied registry status is \texttt{open}. No independent source consultation or literature verification is asserted.
\end{itemize}
\end{document}
